\clearpage

\section{Y and V Phases: Symmetry and Low-Energy
Theory}\label{sec-nbcp-yv-low-energy}

The Y and V reference states carry sublattice-dependent longitudinal and
transverse components. Their common collective coordinate is a
laboratory-\(z\) rotation of all three spins. This chapter identifies
the symmetry and conjugate response shared by the two phases; the
phase-specific spatial and thermal tests are treated separately in
Section~\ref{sec-nbcp-y-phase} and Section~\ref{sec-nbcp-v-phase}.

\subsection{Exact symmetry and accidental classical
degeneracy}\label{exact-symmetry-and-accidental-classical-degeneracy}

At zero bond-dependent SOC, a common spin rotation is an exact U(1)
symmetry of the Hamiltonian. Choosing a transverse orientation then
concerns spontaneous symmetry breaking of that Hamiltonian. At nonzero
SOC, the constant energy of the restricted three-sublattice orbit
instead follows from cancellation of the bond-dependent terms. It is an
accidental classical degeneracy, even though the microscopic Hamiltonian
need not have continuous spin-rotation symmetry. This distinction must
be retained when interpreting a pseudo-Goldstone mode.

\subsection{From a classical orbit to quantum
pinning}\label{from-a-classical-orbit-to-quantum-pinning}

The common spin rotation is
\(\mathbf S_\alpha(\phi)=R_z(\phi)\mathbf S_\alpha(0)\) for all three
sublattices. It is an exact symmetry at \(J_{\rm PD}=J_\Gamma=0\). At
nonzero SOC, the three-sublattice classical energy remains constant
along this orbit by the bond-sum cancellation, although the
finite-momentum spin-wave Hamiltonian changes. The spin-wave vacuum
energy can therefore select discrete values of \(\phi\).

For a stable background and a single soft uniform phase, let
\(C_\phi=\partial_\phi^2e_{\rm sw}|_{\phi_*}\) and
\(\chi_z=\partial m_z/\partial h\), where
\(m_z=(S/3)\sum_\alpha\cos\vartheta_\alpha\) follows the relaxed
classical branch. The leading gap is

\begin{equation}\phantomsection\label{eq-nbcp-gap-summary}{
\Delta_{\rm PG}^2=\frac{C_\phi}{\chi_z}.
}\end{equation}

The conjugate response is a nonuniform polar displacement, not a common
increment of every polar angle. Appendix A derives this statement from
the Berry term; Section~\ref{sec-nbcp-verification-appendix} and its
linked code review retain the legacy comparison. The current
representative gaps at nonzero coupling 0.010 meV are 0.00686217 meV
(Y/PD), \(7.2960\times10^{-6}\) meV (Y/Gamma), 0.0156468 meV (V/PD) and
0.00203898 meV (V/Gamma), for \(B=0.2\) T in Y, \(B=1.4\) T in V and
base mesh \(N=48\). They describe the leading local pinning
calculation at \(T=0\), not a finite-temperature transition or an
all-orders prediction for \(S=1/2\).

\subsection{Order of the angular harmonics in the SOC
couplings}\label{sec-nbcp-yv-selection-rule}

A common spin rotation \(U_\alpha=\exp(-i\alpha\sum_j\hat S_j^z)\) does not
leave the SOC terms invariant, but it maps the Hamiltonian family into
itself. With \(\hat S^\pm\to e^{\mp i\alpha}\hat S^\pm\), the PD bracket
\(\hat S^\pm_i\hat S^\pm_j\) and the Gamma bracket
\(\hat S^\pm_i\hat S^z_j\) acquire phases \(e^{\mp2i\alpha}\) and
\(e^{\mp i\alpha}\). In complex form,
\begin{equation}\label{eq-nbcp-soc-charge}
U_\alpha\hat H(J_{\rm PD},J_\Gamma)U_\alpha^\dagger
=\hat H\bigl(J_{\rm PD}e^{-2i\alpha},J_\Gamma e^{-i\alpha}\bigr),
\end{equation}
where the complex couplings multiply the raising parts and their
conjugates the lowering parts, with the bond phases unchanged. PD
therefore carries spin-rotation charge \(\pm2\) and Gamma charge
\(\pm1\). A rotated background \(\phi\) in the original Hamiltonian is
equivalent to the fixed background \(\phi=0\) in the rotated Hamiltonian.
Any analytic angular energy is then a sum of monomials
\(J_{\rm PD}^{a_+}\bar J_{\rm PD}^{a_-}J_\Gamma^{b_+}\bar J_\Gamma^{b_-}\)
multiplying \(e^{in\phi}\), with
\begin{equation}\label{eq-nbcp-soc-selection}
n=2(a_+-a_-)+(b_+-b_-).
\end{equation}
Real couplings give \(\bar J=J\). The spin-only \(U_\pi\) reverses
\(J_\Gamma\) and shifts \(\phi\) by \(\pi\), so a harmonic with odd
\(n\) is odd in \(J_\Gamma\) and one with even \(n\) is even. This
counting holds at every order of the spin-wave expansion; the lattice
symmetry then restricts which \(n\) survive for each phase.

The lowest monomials follow directly. A sixfold term requires one of
\(J_{\rm PD}^3\), \(J_{\rm PD}^2J_\Gamma^2\), \(J_{\rm PD}J_\Gamma^4\) or
\(J_\Gamma^6\). A threefold term requires \(J_{\rm PD}J_\Gamma\) or
\(J_\Gamma^3\). The existing scans test the first statement. For
pure-Gamma Y, doubling \(J_\Gamma\) from 0.005 to 0.010 meV multiplies
\(\lambda_6\) by 65.7, an effective power of 6.04, which explains the
very small Y/Gamma gap. For pure-PD Y the same doubling gives 12.1, an
effective power of 3.60 instead of 3; the next allowed term,
\(J_{\rm PD}^5\), is therefore not negligible at 0.010 meV, and these
points lie outside the asymptotic weak-SOC regime. The mixed-SOC excess
of the Y sixth harmonic over the sum of the pure-axis values
(Section~\ref{sec-nbcp-y-angular-matching}) is consistent with the
\(J_{\rm PD}^2J_\Gamma^2\) cross term; its separate scaling has not been
fitted.

The saved V scans on the pure axes follow the same counting. Doubling
the coupling multiplies the pure-PD sixth and twelfth harmonics by
effective powers of 3.17 and 6.39 between 0.0025 and 0.005 meV, close to
the leading orders 3 and 6. On the pure-Gamma axis the third, sixth and
ninth harmonics have effective powers 3.01, 6.01 and 9.12.

For V the rule also predicts that near the PD axis the leading threefold
coefficient is \(J_{\rm PD}J_\Gamma\), linear in \(J_\Gamma\), rather than
the \(J_\Gamma^3\) behavior of the pure-Gamma axis. Direct mixed-SOC
calculations confirm this (Section~\ref{sec-nbcp-v-mixed-soc}): at fixed
\(J_{\rm PD}\), \(\lambda_3\) grows linearly in \(J_\Gamma\) and reverses
sign with it, and it exceeds the pure-Gamma value at the same
\(J_\Gamma\) by factors between about \(2\times10^3\) and \(4\times10^5\)
at the sampled points.

\subsection{Canonical response and interpretation of the
gap}\label{canonical-response-and-interpretation-of-the-gap}

The angular curvature alone does not determine the excitation energy.
Its conjugate response allows the polar directions of the three spins to
relax differently, and the susceptibility fixes the normalization of the
collective motion. The derivation in Section~\ref{sec-canonical} and the
Y/V specializations in Section~\ref{sec-states} distinguish this
response from the common polar-angle increment used in the archived
calculation.

The relevant gap is the leading pinning-induced energy of this
collective coordinate. It is not the smallest frequency on a finite
momentum mesh, a gap assigned to an unstable branch, or the energy of a
different internal XXZ mode. A zero-temperature gap also does not by
itself decide whether a thermal algebraic phase exists. The
symmetry-allowed angular harmonics are examined for Y in
Section~\ref{sec-nbcp-y-rg-status} and for V in
Section~\ref{sec-nbcp-v-symmetry}.
